El análisis muestra que la postergación del etiquetado constituye una decisión relevante para la planificación de una viña exportadora, ya que la incertidumbre afecta tanto el volumen total demandado como su distribución entre etiquetas. Frente a ello, la empresa puede diferenciar anticipadamente el producto, reduciendo operaciones posteriores pero comprometiendo inventario, o mantener
botellas sin etiquetar para conservar flexibilidad a cambio de mayores costos de inventario, \textit{set-ups} y utilización de capacidad.

El análisis estadístico verificó la consistencia de la estructura probabilística del árbol de escenarios y evidenció diferencias relevantes entre productos. En particular, las etiquetas del Vino 1 presentan correlaciones predominantemente
negativas, mientras que las del Vino 2 muestran una correlación fuertemente positiva. Esta estructura de dependencia resulta clave para explicar el valor potencial de la postergación, ya que su beneficio no depende únicamente de la variabilidad individual de la demanda, sino también de la posibilidad de reasignar inventario entre productos compatibles.

Asimismo, la holgura observada en el análisis preliminar de capacidad sugiere que, en la instancia base, el valor de la postergación podría provenir en mayor medida de una mejor asignación de inventario y de reducir el costo de decisiones anticipadas incorrectas. No obstante, este resultado es preliminar, pues la utilización efectiva de capacidad dependerá de la interacción entre embotellado, etiquetado, inventarios y \textit{set-ups}. En este contexto, una formulación de optimización estocástica multietapa resulta coherente con la naturaleza del problema, al representar la revelación progresiva de la demanda y permitir decisiones adaptativas. Sin embargo, en esta etapa aún no es posible determinar la política óptima ni cuantificar el beneficio económico de la postergación; dichas relaciones deberán validarse posteriormente mediante el modelo definitivo.. 